\documentclass[11pt,a4paper]{article}
\usepackage[utf8]{inputenc}
\usepackage[T1]{fontenc}
\usepackage{geometry}
\geometry{top=1in, bottom=1in, left=1in, right=1in}
\usepackage{hyperref}
\usepackage{booktabs}
\usepackage{amsmath}
\usepackage{listings}
\usepackage{xcolor}
\usepackage{graphicx}
\usepackage{array}
\usepackage{fancyhdr}
\usepackage{titlesec}
\usepackage{longtable}

\hypersetup{
    colorlinks=true,
    linkcolor=purple,
    filecolor=magenta,      
    urlcolor=cyan,
    pdftitle=Why Genesis?,
    pdfpagemode=FullScreen,
}

\definecolor{codegray}{rgb}{0.96,0.96,0.96}
\definecolor{codeborder}{rgb}{0.85,0.85,0.85}
\definecolor{commentgreen}{rgb}{0.12,0.55,0.22}
\definecolor{keywordblue}{rgb}{0.1,0.2,0.8}

\lstset{
    backgroundcolor=\color{codegray},
    basicstyle=\ttfamily\small,
    breakatwhitespace=false,
    breaklines=true,
    captionpos=b,
    commentstyle=\color{commentgreen},
    frame=single,
    rulecolor=\color{codeborder},
    keepspaces=true,
    keywordstyle=\color{keywordblue},
    language=Python,
    showspaces=false,
    showstringspaces=false,
    showtabs=false,
    tabsize=4
}

\pagestyle{fancy}
\fancyhf{}
\rhead{\small Why Genesis?}
\lhead{\small Project Genesis}
\cfoot{\thepage}

\title{\textbf{Why Genesis?}}
\author{Project Genesis Research Repository}
\date{\small Generated: July 2026}

\begin{document}

\maketitle
\thispagestyle{empty}

\newpage
\section*{Why Genesis?}

Genesis exists to study the emergence of artificial civilizations through reproducible computational experiments. It is an attempt to bridge the gap between complex systems simulation and narrative human history.

\noindent\rule{\textwidth}{0.4pt}

\subsection*{1. The Core Philosophy}

\subsubsection*{Why Simulation?}
Computational simulation is the only laboratory we have to study macro-sociological and evolutionary outcomes that are otherwise impossible to observe in real-time. By simulating agents with cognitive planners, we can observe the emergent friction, cooperation, and collapses that shape complex societies.

\subsubsection*{Why Narrative?}
Data without narrative is noise. If we only expose raw CSV logs and charts, we lose the history. Genesis treats simulations as \textbackslash{}textit{histories} firstusing the \textbackslash{}textbf{Chronicle Engine} to translate numerical events into stories (births, disputes, expansions, declines, collapses). Narrative helps us understand the \textbackslash{}textit{meaning} of the simulation before we analyze the \textbackslash{}textit{mechanisms}.

\subsubsection*{Why Reproducibility?}
Science requires verification. In computational research, this means every published run must be fully reproducible. Genesis treats the \textbackslash{}textbf{Experiment Package} as a first-class, immutable research artifact containing configs, event streams, and raw datasets, enabling anyone to reload, inspect, and verify the experiment locally.

\noindent\rule{\textwidth}{0.4pt}

\subsection*{2. Immutable Truths (What Should Never Change)}

\begin{enumerate}
  \item \textbackslash{}textbf{Decoupled Architecture}: The portal displays and verifies; the simulation engine computes. They interact solely through defined, versioned schemas.
  \item \textbackslash{}textbf{Story Before Charts}: Narrative records rule the Civilization Record. The Observatory (Evidence Explorer) is a secondary research instrument panel.
  \item \textbackslash{}textbf{Reproducibility Over Aesthetics}: Premium design is critical for discovery, but data accuracy and citation governance are absolute.
  \item \textbackslash{}textbf{Emergent Chronicles}: The history of a civilization must emerge organically from events, never scripted or forced.
\end{enumerate}

\end{document}
